%! Polar Function Graphs — Unit 3, Lesson 3.14 — Group Activity (Answer Key)
\documentclass[10pt]{article}
\usepackage{precalculus-article}
\usepackage{precalculus-key}   % pulls in -boxes; do NOT also load -boxes
\newcommand{\TallMath}[1]{$\displaystyle #1\rule[-1.4em]{0pt}{3.2em}$}

\begin{document}

\pageheader{Unit 3, Lesson 3.14}{Group Activity --- Answer Key}
\namepartnerperiod

\vspace{0.08in}

\begin{scenariobox}[Context: Reading Polar Functions]{sky}
In polar coordinates, the function $r = f(\theta)$ assigns a radius to each angle.
As $\theta$ increases, the point $(r, \theta)$ traces a curve in the polar plane.
In this activity your group will read polar functions from tables and graphs,
identify key features, and describe what portions of each curve are traced for
given domains. Choose \textbf{one tier} based on your readiness.
\end{scenariobox}

\vspace{0.08in}

%----------------------------------------------------------------------
% Tier R Key
%----------------------------------------------------------------------
\begin{tcolorbox}[colback=white, colframe=black!40,
    title=\textbf{Tier R --- Reinforce (Key)},
    fonttitle=\bfseries, arc=2mm, left=3mm, right=3mm, top=2mm, bottom=2mm]

\textbf{Directions:} The complete input-output table for $r = 1 + \cos\theta$ is
provided below. Use it to answer the questions.

\vspace{4pt}
\renewcommand{\arraystretch}{1.55}
\begin{tabular}{|c|c|c|c|c|c|c|c|c|c|c|c|c|}
\hline
$\theta$ & $0$ & $\frac{\pi}{6}$ & $\frac{\pi}{3}$ & $\frac{\pi}{2}$ &
           $\frac{2\pi}{3}$ & $\frac{5\pi}{6}$ & $\pi$ &
           $\frac{7\pi}{6}$ & $\frac{4\pi}{3}$ & $\frac{3\pi}{2}$ &
           $\frac{5\pi}{3}$ & $2\pi$ \\
\hline
$r$ & $2$ & $1.87$ & $1.5$ & $1$ & $0.5$ & $0.13$ & $0$ &
      $0.13$ & $0.5$ & $1$ & $1.5$ & $2$ \\
\hline
\end{tabular}

\vspace{8pt}
\textbf{1.}\quad Maximum value of $r$: \ans{2} \quad
At what angle(s)? \ans{$\theta = 0$ (and $\theta = 2\pi$)}

\vspace{6pt}
\textbf{2.}\quad $r = 0$ at $\theta = $ \ans{$\pi$}. \quad
\ans{The polar curve passes through the pole (origin) at that angle.}

\vspace{6pt}
\textbf{3.}\quad On $\theta \in \left[0, \frac{\pi}{2}\right]$, $r$ is
\ans{decreasing}. \ans{The $r$ values in the table go from 2 down to 1 as $\theta$ goes from 0 to $\pi/2$.}

\vspace{6pt}
\textbf{4.}\quad On $\theta \in \left[\frac{\pi}{2}, \pi\right]$, $r$ is
\ans{decreasing} (from 1 to 0).

\vspace{6pt}
\textbf{5.}\quad \ans{Yes. The $r$ values for $\theta$ and $2\pi - \theta$ are equal (e.g., $r(\pi/3) = 1.5 = r(5\pi/3)$), so the curve is symmetric about the $x$-axis.}

\end{tcolorbox}

\vspace{0.08in}

%----------------------------------------------------------------------
% Tier A Key
%----------------------------------------------------------------------
\begin{tcolorbox}[colback=white, colframe=black!40,
    title=\textbf{Tier A --- Apply (Key)},
    fonttitle=\bfseries, arc=2mm, left=3mm, right=3mm, top=2mm, bottom=2mm]

\textbf{Directions:} The rectangular graph of $r = \sin(2\theta)$ for $\theta \in [0, 2\pi]$.

\vspace{4pt}
\begin{center}
\begin{tikzpicture}[x=1.05cm, y=1.15cm]
  \draw[linegray, very thin] (-0.2,-1.25) grid[xstep=0.5, ystep=0.5] (6.7,1.25);
  \draw[->] (-0.2,0) -- (6.9,0) node[right]{\small $\theta$};
  \draw[->] (0,-1.3) -- (0,1.35) node[above]{\small $r$};
  \foreach \x/\lbl in {0.785/$\frac{\pi}{4}$, 1.571/$\frac{\pi}{2}$,
                        2.356/$\frac{3\pi}{4}$, 3.142/$\pi$,
                        3.927/$\frac{5\pi}{4}$, 4.712/$\frac{3\pi}{2}$,
                        5.498/$\frac{7\pi}{4}$, 6.283/$2\pi$} {
    \draw (\x,0.06) -- (\x,-0.06);
    \node[below, font=\scriptsize] at (\x,-0.06) {\lbl};
  }
  \foreach \y/\lbl in {-1/{$-1$}, 1/{$1$}} {
    \draw (0.07,\y) -- (-0.07,\y);
    \node[left, font=\scriptsize] at (-0.07,\y) {\lbl};
  }
  \draw[navy, thick, domain=0:6.3, samples=200]
    plot (\x, {sin(2*deg(\x))});
  \node[right, font=\small, navy] at (6.4, 0.4) {$r = \sin(2\theta)$};
\end{tikzpicture}
\end{center}

\vspace{4pt}
\textbf{1.}\quad \ans{$\theta = 0,\ \pi/2,\ \pi,\ 3\pi/2,\ 2\pi$}

\vspace{6pt}
\textbf{2.}\quad $r > 0$ on: \ans{$(0, \pi/2)$ and $(\pi, 3\pi/2)$}

$r < 0$ on: \ans{$(\pi/2, \pi)$ and $(3\pi/2, 2\pi)$}

\vspace{6pt}
\textbf{3.}\quad \ans{4} arcs. \quad \ans{4} petals.

\vspace{6pt}
\textbf{4.}\quad \ans{When $r < 0$, the point is plotted opposite to $\theta$, so those intervals trace petals in directions $180°$ from the corresponding angle. They still produce petals --- just in different quadrants than you might expect from the angle alone.}

\vspace{6pt}
\textbf{5.}\quad \ans{Yes. The graph shows 4 separate arcs (2 above the axis and 2 below), each zero-to-max/min-to-zero, corresponding to 4 petals on the polar curve.}

\end{tcolorbox}

\vspace{0.08in}

%----------------------------------------------------------------------
% Tier E Key
%----------------------------------------------------------------------
\begin{tcolorbox}[colback=white, colframe=black!40,
    title=\textbf{Tier E --- Extend (Key)},
    fonttitle=\bfseries, arc=2mm, left=3mm, right=3mm, top=2mm, bottom=2mm]

\begin{center}
\begin{tikzpicture}[x=1.3cm, y=1.3cm]
  \foreach \r in {0.5, 1, 1.5, 2} {
    \draw[linegray, very thin] (0,0) circle (\r);
  }
  \foreach \ang in {0, 30, 60, 90, 120, 150, 180, 210, 240, 270, 300, 330} {
    \draw[linegray, very thin] (0,0) -- ({2.2*cos(\ang)}, {2.2*sin(\ang)});
  }
  \draw[->] (-2.3,0) -- (2.4,0) node[right]{\small $x$};
  \draw[->] (0,-2.3) -- (0,2.4) node[above]{\small $y$};
  \draw[navy, very thick, domain=0:3.1416, samples=200, smooth]
    plot ({2*sin(deg(\x))*cos(\x r)}, {2*sin(deg(\x))*sin(\x r)});
  \draw[redacc, very thick, domain=0:6.2832, samples=300, smooth]
    plot ({(1+sin(deg(\x)))*cos(\x r)}, {(1+sin(deg(\x)))*sin(\x r)});
  \node[font=\small, navy] at (0.8, 1.9) {$r = 2\sin\theta$};
  \node[font=\small, redacc] at (1.7, 0.9) {$r = 1+\sin\theta$};
  \node[below right, font=\scriptsize] at (1,0) {$1$};
  \node[below right, font=\scriptsize] at (2,0) {$2$};
  \node[above right, font=\scriptsize] at (0,1) {$1$};
  \node[above right, font=\scriptsize] at (0,2) {$2$};
\end{tikzpicture}
\end{center}

\vspace{4pt}
\textbf{1.}\quad $r = 2\sin\theta$:
Maximum $r$: \ans{2} \quad
$r = 0$ at: \ans{$\theta = 0$ and $\theta = \pi$}

\vspace{4pt}
\textbf{2.}\quad $r = 1 + \sin\theta$:
Maximum $r$: \ans{2} at $\theta = $ \ans{$\pi/2$} \quad
Minimum $r$: \ans{0} at $\theta = $ \ans{$3\pi/2$} \quad
$r = 0$ at: \ans{$\theta = 3\pi/2$}

\vspace{6pt}
\textbf{3.}\quad \ans{For $\theta \in [0, \pi]$, $r = 2\sin\theta \geq 0$: the curve starts at the pole, moves away (maximum $r = 2$ at $\theta = \pi/2$), and returns to the pole. For $\theta \in (\pi, 2\pi)$, $r < 0$, so the points are plotted in the opposite direction --- they retrace the same circle. The full circle is actually completed in $\theta \in [0, \pi]$.}

\vspace{6pt}
\textbf{4.}\quad \ans{$r = 1 + \sin\theta$ is a cardioid (heart-shaped), while $r = 2\sin\theta$ is a circle. The cardioid has maximum $r = 2$ (same as the circle) but passes through the pole only once (at $\theta = 3\pi/2$). The cardioid extends into all four quadrants, while the circle sits entirely in the upper half-plane.}

\end{tcolorbox}

\end{document}
